\documentclass[aspectratio=169,12pt,fontset=none]{ctexbeamer}
% XeLaTeX typography for less clear presentation screens.
% Font files and redistribution licenses are bundled under fonts/.
\usepackage{fontspec}
\setsansfont[
  Path=fonts/,
  BoldFont=OpenSans-Bold.ttf,
  ItalicFont=OpenSans-SemiboldItalic.ttf,
  BoldItalicFont=OpenSans-BoldItalic.ttf
]{OpenSans-Semibold.ttf}
\ifdefined\setCJKmainfont
  \setCJKmainfont[Path=fonts/,BoldFont=NotoSansSC-Bold.otf]{NotoSansSC-Medium.otf}
  \setCJKsansfont[Path=fonts/,BoldFont=NotoSansSC-Bold.otf]{NotoSansSC-Medium.otf}
  \setCJKmonofont[Path=fonts/,BoldFont=NotoSansSC-Bold.otf]{NotoSansSC-Medium.otf}
\fi

\usepackage{amsmath}
\usepackage{pgfplots}
\pgfplotsset{compat=1.18}
\usetheme{Arclight}
\arclightsetcompany{弧光量子}
\usetikzlibrary{arrows.meta,positioning}

\title[量子校准与数据解读]{从信号到判断}
\subtitle{量子校准与数据解读 / 公司报告样式预览}
\author{报告人：待填写}
\institute{弧光量子}
\date{2026 年 10 月 12 日}

\begin{document}
\begin{frame}[plain]
  \titlepage
\end{frame}
\arclightagenda{报告目录}

\section{从实验到数据}
\arclightdivider{01\quad 从实验到数据}{先明确：我们控制什么，测量什么？}

\begin{frame}{控制信号通过实验变成测量记录}
  \vspace{4mm}
  \begin{figure}
 \centering
    \begin{tikzpicture}[
      node distance=11mm,
      stage/.style={draw=ArclightRule,line width=.7pt,fill=ArclightPale,
        minimum width=34mm,minimum height=21mm,align=center},
      flow/.style={-{Latex[length=2.5mm]},draw=ArclightRed,line width=1pt}]
      \node[stage] (pulse) {控制信号\\[2mm]{\small Drive pulse}};
      \node[stage,right=of pulse] (device) {器件响应\\[2mm]{\small Qubit + resonator}};
      \node[stage,right=of device] (record) {测量记录\\[2mm]{\small Digitized IQ}};
      \draw[flow] (pulse) -- (device);
      \draw[flow] (device) -- (record);
    \end{tikzpicture}
  \end{figure}
  \vspace{5mm}
  \textbf{校准建立控制参数与测量响应之间的关系。}
  \vfill
  {\footnotesize\color{ArclightMuted}概念示意；完整解释见配套笔记。}
\end{frame}

\begin{frame}{平坦曲线无法确定衰减时间}
  \begin{columns}[T,onlytextwidth]
    \begin{column}{.47\textwidth}
      \textbf{窗口内可见衰减}\par
      \vspace{2mm}
      \begingroup
      \setlength{\topsep}{0pt}\setlength{\partopsep}{0pt}
      \begin{figure}
      \centering
      \begin{tikzpicture}
        \begin{axis}[width=59mm,height=36mm,axis lines=left,
          xmin=0,xmax=6.4,ymin=0,ymax=1.18,ticks=none,
          xlabel={$t$},ylabel={$y$},clip=false]
          \addplot[ArclightInk,line width=1.1pt,domain=0:6,samples=60] {exp(-x/2)};
          \node[font=\small,align=center] at (axis cs:3.2,.77) {曲线发生\\明显变化};
        \end{axis}
      \end{tikzpicture}
      \end{figure}
      \endgroup
    \end{column}
    \begin{column}{.47\textwidth}
      \textbf{窗口内几乎平坦}\par
      \vspace{2mm}
      \begingroup
      \setlength{\topsep}{0pt}\setlength{\partopsep}{0pt}
      \begin{figure}
      \centering
      \begin{tikzpicture}
        \begin{axis}[width=59mm,height=36mm,axis lines=left,
          xmin=0,xmax=6.4,ymin=0,ymax=1.18,ticks=none,
          xlabel={$t$},ylabel={$y$},clip=false]
          \addplot[ArclightInk,line width=1.1pt,domain=0:6,samples=60] {.96-.002*x};
          \draw[ArclightRed,line width=.8pt,-{Latex}] (axis cs:3,.60) -- (axis cs:4,.94);
          \node[font=\small,align=center] at (axis cs:2.9,.44) {没有可分辨\\的衰减};
        \end{axis}
      \end{tikzpicture}
      \end{figure}
      \endgroup
    \end{column}
  \end{columns}
  \vspace{3mm}
  \textbf{拟合返回的数值，需要曲线提供支持。}
  \vfill
  {\footnotesize\color{ArclightMuted}示意曲线，非实验数据。}
\end{frame}

\section{交流与讨论}
\arclightdivider{交流与讨论}{问题、观察与下一步}
\end{document}
